\chapter{Lesson 11: Sports, and a Review of Complements}

\section{Vocabulary}

\begin{center}
\begin{tabular}{llll}
\toprule
Characters & Pinyin & Part of speech & Meaning \\
\midrule
\zh{运动} & yùndòng & noun / verb & sports; to exercise \\
\zh{锻炼} & duànliàn & verb & to work out \\
\zh{身体} & shēntǐ & noun & body, health \\
\zh{跑步} & pǎo bù & verb-object & to jog \\
\zh{游泳} & yóu yǒng & verb-object & to swim \\
\zh{打篮球} & dǎ lánqiú & verb-object & to play basketball \\
\zh{打网球} & dǎ wǎngqiú & verb-object & to play tennis \\
\zh{打乒乓球} & dǎ pīngpāngqiú & verb-object & to play ping-pong \\
\zh{踢足球} & tī zúqiú & verb-object & to play soccer \\
\zh{打冰球} & dǎ bīngqiú & verb-object & to play hockey \\
\zh{滑冰} & huá bīng & verb-object & to skate \\
\zh{滑雪} & huá xuě & verb-object & to ski \\
\zh{比赛} & bǐsài & noun / verb & match; to compete \\
\zh{赢} & yíng & verb & to win \\
\zh{输} & shū & verb & to lose \\
\zh{球队} & qiúduì & noun & team \\
\zh{球迷} & qiúmí & noun & fan \\
\zh{为了} & wèile & preposition & in order to \\
\zh{一边\ldots{}一边} & yìbiān \ldots yìbiān & & while (doing two things) \\
\zh{好处} & hǎochu & noun & benefit \\
\zh{被} & bèi & preposition & (passive marker) \\
\zh{小时} & xiǎoshí & noun & hour \\
\zh{半天} & bàntiān & noun & a long time \\
\bottomrule
\end{tabular}
\end{center}

\section{Verbs for sports}

\begin{important}[\zh{打}, \zh{踢} and the rest]
Ball games played with the hands use \zh{打}: \zh{打篮球, 打网球, 打排球, 打乒乓球},
and also \zh{打冰球} (with a stick). Soccer uses \zh{踢} ``kick''. Swimming, running,
skating and skiing are verb--object compounds: \zh{游泳, 跑步, 滑冰, 滑雪}.
\end{important}

\section{Grammar}

\subsection{Duration of an action}

The duration phrase goes \emph{after} the verb (unlike time-when, which goes before).
\begin{itemize}
  \item \zh{我每天跑步跑半个小时。} I run for half an hour every day.
  \item \zh{他游了一个小时(的)泳。} He swam for an hour.
  \item \zh{我学中文学了两年了。} I've been studying Chinese for two years (and still am).
\end{itemize}

\begin{important}[Duration with an object]
Either repeat the verb (\zh{打球打了两个小时}) or put the duration between the verb and
the object, optionally with \zh{的} (\zh{打了两个小时(的)球}). A personal pronoun object
comes before the duration: \zh{我等了你半天}.
\end{important}

\begin{center}
\begin{tikzpicture}[x=1cm]
  \draw[thick, -{Stealth}] (0,0) -- (7,0) node[right, font=\small] {time};
  \fill[purple] (2,0) circle (0.12);
  \node[above, font=\small, purple, align=center] at (2,0.2) {time-when\\ \zh{八点}};
  \draw[line width=4pt, teal] (2,-0.5) -- (6,-0.5);
  \node[below, font=\small, teal, align=center] at (4,-0.6) {duration \zh{两个小时}};
  \node[font=\small, anchor=west] at (8.2,0.8) {\zh{我\textbf{八点}去打球，}};
  \node[font=\small, anchor=west] at (8.2,0.2) {\zh{打了\textbf{两个小时}。}};
  \node[font=\scriptsize, anchor=west, align=left] at (8.2,-0.7) {when: before the verb\\ how long: after the verb};
\end{tikzpicture}
\end{center}

\subsection{\zh{为了} ``in order to''}

\zh{为了} + purpose comes at the start of the sentence.
\begin{itemize}
  \item \zh{为了锻炼身体，他每天早上去游泳。} To keep fit, he swims every morning.
  \item \zh{为了看这场比赛，我们排了三个小时的队。} We queued for three hours to see the
        game.
\end{itemize}

\subsection{\zh{一边\ldots{}一边\ldots}}

Two actions done at the same time by the same person:
\begin{itemize}
  \item \zh{他一边跑步，一边听音乐。} He listens to music while running.
  \item \zh{我们一边看冰球比赛，一边吃比萨饼。} We ate pizza while watching the hockey game.
\end{itemize}

\subsection{The passive with \zh{被} (introduction)}

\zh{被} marks the doer in a passive sentence; traditionally it describes something
unfortunate.
\begin{itemize}
  \item \zh{我的自行车被人偷了。} My bike was stolen.
  \item \zh{我们队被他们打败了。} Our team was beaten by theirs.
\end{itemize}
Like \zh{把}, the verb needs something after it (\zh{了}, a complement). In speech
\zh{让} and \zh{叫} can replace \zh{被} when the doer is mentioned.

\section{Review: the complements}

The final exam asked us to identify and use every type of complement learned in 101R
and 102R. All of them follow the verb. A potential complement is a result or direction
complement with \zh{得} or \zh{不} inserted.

\begin{center}
\begin{tikzpicture}[h/.style={draw, very thick, rounded corners, fill=yellow!25, font=\bfseries}, c/.style={draw, rounded corners, align=center, font=\small, fill=#1, text width=3.6cm}]
  \node[h] (V) at (0,0) {Verb + complement};
  \node[c=blue!8] (deg) at (-4.6,2.4) {degree\\ V\zh{得} + adj\\ \zh{跑得很快}};
  \node[c=green!10] (res) at (0,3.0) {result\\ V + adj/V\\ \zh{吃完, 听懂, 找到}};
  \node[c=orange!12] (dir) at (4.6,2.4) {direction\\ V + \zh{来/去}, \zh{上来}\ldots\\ \zh{走进来, 拿出去}};
  \node[c=red!8] (pot) at (-4.6,-2.4) {potential\\ V\zh{得/不}C\\ \zh{听得懂, 买不到}};
  \node[c=purple!8] (dur) at (0,-3.0) {duration\\ V + time span\\ \zh{游了一个小时}};
  \node[c=gray!12] (frq) at (4.6,-2.4) {frequency\\ V + number \zh{次}\\ \zh{去过两次}};
  \foreach \n in {deg,res,dir,pot,dur,frq} \draw[thick, gray] (V) -- (\n);
\end{tikzpicture}
\end{center}

\begin{center}
\begin{tabular}{llll}
\toprule
Type & Negative & Question & Object position \\
\midrule
degree & V\zh{得不}adj & V\zh{得怎么样？} & repeat verb or front object \\
result & \zh{没}VC & VC\zh{了没有？} & after complement \\
direction & \zh{没}VC & VC\zh{了吗？} & place before \zh{来/去} \\
potential & V\zh{不}C & V\zh{得}C V\zh{不}C? & after or topicalized \\
duration & \zh{没}V + duration (rare) & V\zh{了多长时间？} & repeat verb or between \\
frequency & \zh{没}V\zh{过} & V\zh{过几次？} & pronoun before \zh{次} \\
\bottomrule
\end{tabular}
\end{center}

\begin{example}[Identify the complement]\leavevmode
\begin{enumerate}
  \item \zh{他冰球打得非常好。} --- degree.
  \item \zh{老师说的话你都听清楚了吗？} --- result.
  \item \zh{比赛票已经卖完了，买不到了。} --- result (\zh{卖完}), potential (\zh{买不到}).
  \item \zh{球迷们都站起来了。} --- direction.
  \item \zh{我们练习了三个星期。} --- duration.
\end{enumerate}
\end{example}

\section{Culture: sports}

\begin{aside}
\begin{itemize}
  \item Ping-pong is called China's ``national ball'' (\zh{国球}). The 1971 ``ping-pong
        diplomacy'' (\zh{乒乓外交}) visit of the US team helped open US--China relations.
  \item Beijing hosted the 2008 Summer Olympics and the 2022 Winter Olympics, the first
        city to host both.
  \item Morning exercise in parks: \zh{太极拳} tàijíquán and square dancing
        (\zh{广场舞} guǎngchǎngwǔ, mostly retirees).
  \item Hockey (\zh{冰球}) is ``Canada's game''; the Chinese name for the NHL is
        \zh{国家冰球联盟}.
\end{itemize}
Traditional forms: \zht{運動} sports, \zht{籃球} basketball, \zht{體育} physical education.
\end{aside}

\section{Model composition}

\zh{我最喜欢的运动是游泳。为了锻炼身体，我每个星期去学校的游泳池游三次泳，每次游四十分钟。游泳的好处很多：不但对身体好，而且游完以后心情也越来越好。冬天滑铁卢太冷了，我不能在外面跑步，所以我也常常跟朋友一起去滑冰。我滑得不太好，可是我们一边滑冰一边聊天，很有意思。上个星期我们去看了一场冰球比赛，我们的球队输了，大家都很失望。}

(\zh{游泳池} yóuyǒngchí = pool; \zh{心情} xīnqíng = mood; \zh{场} chǎng = MW for games;
\zh{失望} shīwàng = disappointed.)

\section*{Review}
\begin{itemize}
  \item \zh{打} for hand-ball games and hockey, \zh{踢} for soccer.
  \item Duration after the verb; repeat the verb with an object; V\zh{了}+duration+\zh{了}
        for ongoing.
  \item \zh{为了} purpose first; \zh{一边\ldots{}一边\ldots} simultaneous actions.
  \item \zh{被} passive (adversity), needs something after the verb.
  \item Six complements: degree, result, direction, potential, duration, frequency.
\end{itemize}
